\documentclass[11pt,a4paper]{article}

\usepackage[utf8]{inputenc}
\usepackage[T1]{fontenc}
\usepackage[french]{babel}
\usepackage{amsmath,amssymb,amsfonts}
\usepackage{geometry}
\usepackage{graphicx}
\usepackage{booktabs}
\usepackage{hyperref}
\usepackage{array}
\usepackage{xcolor}
\usepackage{float}

\geometry{margin=2.5cm}
\graphicspath{{figures/}}

\title{
\textbf{Limite spectrale continue du Laplacien d'intrication BuP}\\
\large Des valeurs propres discrètes au Laplacien de Laplace--Beltrami
}

\author{Farid Hamdad}
\date{\2026}

\begin{document}

\maketitle

\begin{abstract}
Paper~15 a fourni une première validation numérique de la limite géométrique
de Bottom-Up Quantum Gravity en montrant que le Laplacien d'intrication
rescalé
\[
L_\epsilon=\frac{D-W}{\epsilon}
\]
reconstruit la dimension spectrale de géométries contrôlées. Le présent papier
renforce cette validation en passant d'un test global de trace de chaleur à
un test spectral mode par mode. Nous comparons les premières valeurs propres
du Laplacien BuP rescalé aux spectres analytiques du cercle \(S^1\) et du tore
plat \(T^2\). Sur \(S^1\), les 12 premiers modes non nuls sont reconstruits
avec une erreur relative moyenne de \(0.0051\) à \(N=1024\). Sur \(T^2\), les
20 premiers modes non nuls sont reconstruits avec une erreur relative moyenne
de \(0.0266\) à \(N=1024\). Ces résultats soutiennent la convergence
spectrale bas-modes
\[
c_N L_N
=
c_N\frac{D_N-W_N}{\epsilon_N}
\longrightarrow
-\Delta_g,
\]
où \(c_N\) est une normalisation scalaire encore ajustée numériquement. Le
problème central restant est de dériver analytiquement \(c_N\) et de préciser
le mode de convergence : convergence ponctuelle sur fonctions tests,
convergence de formes quadratiques, convergence de trace de chaleur ou
convergence spectrale.
\end{abstract}

\tableofcontents

\section{Introduction}

Le programme Bottom-Up Quantum Gravity propose que la géométrie ne soit pas
fondamentale, mais émerge de la structure d'intrication d'un état quantique
global. L'objet de départ n'est donc pas une métrique \(g_{\mu\nu}\), mais une
matrice relationnelle
\[
W_{ij}=I(i:j),
\]
où \(I(i:j)\) désigne l'information mutuelle entre les degrés de liberté
\(i\) et \(j\).

À partir de \(W_{ij}\), on construit un Laplacien d'intrication :
\[
L_{\rm ent}=D-W,
\qquad
D_{ii}=\sum_j W_{ij}.
\]
Ce Laplacien encode la connectivité, la diffusion et la structure spectrale du
graphe d'intrication.

Paper~15 a montré numériquement que le Laplacien rescalé
\[
L_\epsilon=\frac{D-W}{\epsilon}
\]
récupère la dimension spectrale attendue sur plusieurs géométries contrôlées.
La dimension spectrale était extraite de la trace de chaleur :
\[
Z(t)=\mathrm{Tr}\,e^{-tL},
\]
par
\[
d_s(t)
=
-2\frac{d\log Z(t)}{d\log t}.
\]

Le présent papier va plus loin. Au lieu de tester uniquement une quantité
globale dérivée de la trace de chaleur, nous testons directement la convergence
des valeurs propres individuelles :
\[
\lambda_k(L_N)
\longrightarrow
\lambda_k(\Delta_g).
\]

La question centrale de Paper~16 est donc :
\[
\boxed{
\text{Le Laplacien BuP rescalé converge-t-il spectralement vers le
Laplacien de Laplace--Beltrami ?}
}
\]

La forme recherchée est :
\[
c_N L_N
=
c_N\frac{D_N-W_N}{\epsilon_N}
\longrightarrow
-\Delta_g,
\]
où \(c_N\) est une constante de normalisation dépendant de la discrétisation,
du noyau, de la dimension et de l'échelle \(\epsilon_N\).

\section{Laplacien d'intrication BuP}

Le point de départ est la matrice d'information mutuelle :
\[
W_{ij}=I(i:j).
\]

On définit le degré pondéré :
\[
D_{ii}=\sum_j W_{ij},
\]
puis le Laplacien discret non normalisé :
\[
L_N^0=D_N-W_N.
\]

Le test de Paper~15 a montré que le Laplacien normalisé standard
\[
L_{\rm norm}=I-D^{-1/2}WD^{-1/2}
\]
n'est pas le bon opérateur pour approcher directement un Laplacien continu.
Il tend à donner une dimension spectrale artificielle, notamment en dimension
deux.

Le bon objet numérique identifié dans Paper~15 est le Laplacien rescalé :
\[
L_N=\frac{D_N-W_N}{\epsilon_N}.
\]

L'échelle \(\epsilon_N\) est obtenue à partir de la distance typique au
voisinage local. Dans les tests numériques, on utilise :
\[
\epsilon_N
=
\gamma\,\epsilon_{\rm knn},
\]
avec \(\gamma=0.5\), et \(k=\sqrt{N}\). Le noyau utilisé est de type gaussien :
\[
W_{ij}
=
\exp\left[
-\frac{d(i,j)^2}{4\epsilon_N}
\right],
\]
restreint à un graphe local \(k\)-NN.

La cible continue est le Laplacien de Laplace--Beltrami :
\[
\Delta_g
=
\frac{1}{\sqrt{|g|}}
\partial_\mu
\left(
\sqrt{|g|}g^{\mu\nu}\partial_\nu
\right).
\]

En convention positive, le spectre de \(-\Delta_g\) est non négatif :
\[
-\Delta_g\phi_k=\lambda_k\phi_k,
\qquad
\lambda_k\geq0.
\]

Ainsi, la convergence attendue est :
\[
c_N L_N\to -\Delta_g.
\]

\section{Héritage de Paper 15 : convergence de dimension spectrale}

Paper~15 a établi une première validation numérique de la limite
\(L_N\to\Delta_g\) au niveau de la dimension spectrale.

Avec le Laplacien rescalé
\[
L_\epsilon=\frac{D-W}{\epsilon},
\]
les dimensions spectrales mesurées étaient :

\[
\begin{array}{lccc}
\toprule
{\rm Géométrie} & d_{\rm cible} & d_s^{\rm mesuré} & {\rm erreur} \\
\midrule
{\rm cercle} & 1 & 1.0021 & 0.0021 \\
{\rm intervalle} & 1 & 0.9445 & 0.0555 \\
{\rm grille\ 2D} & 2 & 1.9938 & 0.0062 \\
{\rm sphère} & 2 & 1.8964 & 0.1036 \\
\bottomrule
\end{array}
\]

Ces résultats soutenaient la flèche :
\[
L_N\to\Delta_g
\]
au niveau de la trace de chaleur.

Cependant, la dimension spectrale est une quantité globale. Elle peut être
correcte même si les valeurs propres individuelles présentent encore des
écarts. Paper~16 cherche donc à tester une condition plus forte :
\[
\lambda_k(c_NL_N)
\simeq
\lambda_k(-\Delta_g).
\]

\section{Problème de Paper 16 : convergence spectrale bas-modes}

La convergence spectrale d'un opérateur discret vers un opérateur continu peut
être testée en comparant leurs premières valeurs propres. Pour un opérateur
elliptique compact, les bas modes encodent la géométrie à grande échelle.

Le problème étudié ici est :
\[
\lambda_k(c_NL_N)
\to
\lambda_k(-\Delta_g),
\qquad
k=1,\ldots,K,
\]
pour des géométries dont le spectre analytique est connu.

Nous utilisons deux cas tests :

\begin{enumerate}
    \item le cercle unité \(S^1\), de dimension \(1\) ;
    \item le tore plat unité \(T^2=[0,1)^2\), de dimension \(2\).
\end{enumerate}

Ces deux géométries sont particulièrement utiles parce qu'elles sont compactes,
sans bord, et possèdent un spectre analytique explicite. Elles permettent donc
de tester la convergence spectrale sans ambiguïté due aux conditions de bord.

\section{Test spectral sur le cercle \(S^1\)}

\subsection{Spectre analytique}

Sur le cercle unité \(S^1\), paramétré par
\[
\theta\in[0,2\pi),
\]
le Laplacien positif possède le spectre :
\[
\lambda_m=m^2,
\qquad
m=0,1,2,\ldots
\]

Les modes \(m\geq1\) sont dégénérés, correspondant aux fonctions
\[
\cos(m\theta),
\qquad
\sin(m\theta).
\]

Ainsi, le spectre non nul commence par :
\[
1,1,4,4,9,9,16,16,\ldots
\]

\subsection{Construction du graphe}

Nous plaçons \(N\) points uniformément sur le cercle :
\[
\theta_i=\frac{2\pi i}{N},
\qquad
i=0,\ldots,N-1.
\]

La distance géodésique intrinsèque est :
\[
d(\theta_i,\theta_j)
=
\min\left(
|\theta_i-\theta_j|,
2\pi-|\theta_i-\theta_j|
\right).
\]

Les poids sont définis par :
\[
W_{ij}
=
\exp\left[
-\frac{d(\theta_i,\theta_j)^2}{4\epsilon_N}
\right],
\]
sur un graphe local \(k\)-NN avec :
\[
k=\sqrt{N}.
\]

Le Laplacien testé est :
\[
L_N=\frac{D_N-W_N}{\epsilon_N}.
\]

Comme la normalisation exacte du noyau n'est pas encore dérivée, on ajuste un
facteur scalaire \(c_N\) tel que :
\[
c_N\lambda_k^{\rm graph}
\simeq
\lambda_k^{S^1}.
\]

Le facteur \(c_N\) est obtenu par moindres carrés sur les premiers modes :
\[
c_N
=
\frac{
\sum_{k=1}^{K_{\rm fit}}
\lambda_k^{\rm graph}\lambda_k^{\rm target}
}{
\sum_{k=1}^{K_{\rm fit}}
(\lambda_k^{\rm graph})^2
}.
\]

\subsection{Résultats}

Les erreurs relatives sur les 12 premiers modes non nuls décroissent
systématiquement avec \(N\) :

\[
\begin{array}{cccc}
\toprule
N & {\rm erreur\ moyenne} & {\rm erreur\ médiane} & {\rm erreur\ maximale}\\
\midrule
64 & 0.0937 & 0.0952 & 0.1852\\
128 & 0.0492 & 0.0480 & 0.1003\\
256 & 0.0214 & 0.0215 & 0.0446\\
512 & 0.0109 & 0.0094 & 0.0249\\
1024 & 0.0051 & 0.0051 & 0.0107\\
\bottomrule
\end{array}
\]

À \(N=1024\), le premier mode non nul est reconstruit comme :
\[
\lambda_1^{\rm scaled}=1.0062,
\]
à comparer à la valeur analytique :
\[
\lambda_1^{S^1}=1.
\]

L'erreur relative sur ce premier mode est donc :
\[
\frac{|1.0062-1|}{1}
\simeq
0.0062.
\]

Cela fournit une évidence numérique forte pour :
\[
c_NL_N
\longrightarrow
-\Delta_{S^1}.
\]

\section{Test spectral sur le tore plat \(T^2\)}

\subsection{Spectre analytique}

Le tore plat unité est :
\[
T^2=[0,1)^2,
\]
avec conditions périodiques.

Les modes propres sont :
\[
e^{2\pi i(mx+ny)},
\qquad
(m,n)\in\mathbb{Z}^2.
\]

Le spectre positif de \(-\Delta_{T^2}\) est :
\[
\lambda_{m,n}
=
4\pi^2(m^2+n^2),
\qquad
(m,n)\neq(0,0).
\]

Le premier mode non nul est donc :
\[
\lambda_1^{T^2}=4\pi^2\simeq39.4784.
\]

\subsection{Construction du graphe}

Nous utilisons une grille périodique \(m\times m\), avec :
\[
N=m^2.
\]

La distance intrinsèque périodique entre deux points
\[
x_i=(u_i,v_i),
\qquad
x_j=(u_j,v_j)
\]
est :
\[
d(x_i,x_j)^2
=
\min(|u_i-u_j|,1-|u_i-u_j|)^2
+
\min(|v_i-v_j|,1-|v_i-v_j|)^2.
\]

Les poids sont :
\[
W_{ij}
=
\exp\left[
-\frac{d(x_i,x_j)^2}{4\epsilon_N}
\right],
\]
sur un graphe local \(k\)-NN.

Le Laplacien testé est :
\[
L_N=\frac{D_N-W_N}{\epsilon_N}.
\]

Comme pour le cercle, un facteur \(c_N\) est ajusté :
\[
c_N\lambda_k^{\rm graph}
\simeq
\lambda_k^{T^2}.
\]

\subsection{Résultats}

Les erreurs relatives sur les 20 premiers modes non nuls décroissent avec la
résolution :

\[
\begin{array}{cccc}
\toprule
N & {\rm erreur\ moyenne} & {\rm erreur\ médiane} & {\rm erreur\ maximale}\\
\midrule
64 & 0.0951 & 0.1340 & 0.1467\\
144 & 0.0757 & 0.0437 & 0.1367\\
256 & 0.0604 & 0.0574 & 0.1850\\
576 & 0.0334 & 0.0455 & 0.0504\\
1024 & 0.0266 & 0.0263 & 0.0522\\
\bottomrule
\end{array}
\]

À \(N=1024\), on obtient :
\[
\lambda_1^{\rm scaled}=40.9613.
\]

La valeur analytique est :
\[
\lambda_1^{T^2}
=
4\pi^2
=
39.4784.
\]

L'erreur relative du premier mode est :
\[
\frac{|40.9613-39.4784|}{39.4784}
\simeq
0.0376.
\]

Le tore plat présente une convergence moins rapide que le cercle, ce qui est
attendu : en dimension deux, les dégénérescences spectrales sont plus
nombreuses, les sous-espaces propres sont plus sensibles à la discrétisation,
et le graphe \(k\)-NN peut légèrement briser les symétries exactes du tore.

Néanmoins, l'erreur moyenne finale sur les 20 premiers modes est :
\[
0.0266,
\]
ce qui fournit une évidence numérique positive pour :
\[
c_NL_N
\longrightarrow
-\Delta_{T^2}.
\]

\section{Résumé des résultats spectraux}

Les deux tests principaux de Paper~16 sont résumés dans le tableau suivant :

\[
\begin{array}{llllllll}
\toprule
{\rm Test} & {\rm Limite} & N & {\rm modes} &
{\rm erreur\ moy.} & {\rm erreur\ méd.} &
{\rm erreur\ max.} & \lambda_1^{\rm scaled} \\
\midrule
S^1 &
c_NL_N\to-\Delta_{S^1} &
1024 & 12 & 0.0051 & 0.0051 & 0.0107 & 1.0062 \\
T^2 &
c_NL_N\to-\Delta_{T^2} &
1024 & 20 & 0.0266 & 0.0263 & 0.0522 & 40.9613 \\
\bottomrule
\end{array}
\]

Pour comparaison, les valeurs cibles du premier mode sont :
\[
\lambda_1^{S^1}=1,
\qquad
\lambda_1^{T^2}=4\pi^2=39.4784.
\]

Ainsi, Paper~16 renforce Paper~15 :

\[
\text{Paper 15 : }
d_s\to D,
\]

\[
\text{Paper 16 : }
\lambda_k(c_NL_N)\to\lambda_k(-\Delta_g).
\]

\section{Interprétation}

Ces résultats montrent que le Laplacien BuP rescalé ne reproduit pas seulement
une dimension spectrale globale. Il reconstruit aussi le bas spectre du
Laplacien de Laplace--Beltrami sur des géométries contrôlées.

Le passage de Paper~15 à Paper~16 est important :

\[
\mathrm{Tr}\,e^{-tL_N}
\quad
\text{teste une propriété globale},
\]
tandis que
\[
\lambda_k(L_N)
\quad
\text{teste les modes individuels}.
\]

La convergence des valeurs propres est donc une validation plus fine de la
limite spectrale.

Le résultat actuel peut être formulé ainsi :
\[
\boxed{
\text{Le bas spectre du Laplacien BuP rescalé converge numériquement vers le
bas spectre de }-\Delta_g\text{ sur }S^1\text{ et }T^2.
}
\]

Il reste toutefois une étape théorique majeure : la normalisation \(c_N\).

\section{Le problème de la normalisation \(c_N\)}

Dans les tests actuels, \(c_N\) est ajusté numériquement. Ce facteur absorbe
plusieurs effets :

\begin{itemize}
    \item la normalisation du noyau gaussien ;
    \item la densité d'échantillonnage ;
    \item la dimension de la variété ;
    \item le choix du graphe \(k\)-NN ;
    \item le facteur \(\epsilon_N\) ;
    \item la différence entre une intégrale continue et une somme discrète.
\end{itemize}

Le problème central de Paper~16 est donc de remplacer l'ajustement :
\[
c_N
=
\operatorname*{argmin}_c
\sum_k
\left(
c\lambda_k^{\rm graph}
-
\lambda_k^{\rm target}
\right)^2
\]
par une formule analytique :
\[
c_N=C(D,\rho_N,\epsilon_N,k_N,\gamma),
\]
où \(D\) est la dimension, \(\rho_N\) la densité d'échantillonnage, et
\(\gamma\) le facteur définissant \(\epsilon_N\).

Un résultat théorique satisfaisant aurait la forme :
\[
c_N\frac{D_N-W_N}{\epsilon_N}
\longrightarrow
-\Delta_g.
\]

\section{Vers un théorème de convergence}

Un énoncé cible possible est le suivant.

\medskip

\noindent
\textbf{Énoncé cible.}
Soit \(M\) une variété riemannienne compacte de dimension \(D\), échantillonnée
par des points \(x_1,\ldots,x_N\). Supposons que les poids BuP vérifient
localement :
\[
W_{ij}
=
\exp\left[
-\frac{d_g(x_i,x_j)^2}{4\epsilon_N}
\right]
+
o(1),
\]
avec :
\[
\epsilon_N\to0,
\qquad
N\epsilon_N^{D/2}\to\infty.
\]

Alors il existe une normalisation \(c_N\) telle que :
\[
c_N\frac{D_N-W_N}{\epsilon_N}
\longrightarrow
-\Delta_g,
\]
dans un mode de convergence à préciser.

\medskip

Les modes de convergence possibles sont :

\begin{enumerate}
    \item convergence ponctuelle sur fonctions tests lisses ;
    \item convergence des formes quadratiques ;
    \item convergence de Mosco ;
    \item convergence des valeurs propres ;
    \item convergence des noyaux de chaleur ;
    \item convergence de la trace de chaleur ;
    \item convergence de la dimension spectrale.
\end{enumerate}

Dans le cadre BuP, la difficulté supplémentaire est que les poids ne sont pas
initialement posés comme un noyau géométrique. Ils proviennent de
l'intrication :
\[
W_{ij}=I(i:j).
\]
Il faut donc établir sous quelles conditions \(I(i:j)\) se comporte
localement comme un noyau thermique ou gaussien.

\section{Lien avec l'action spectrale BuP}

Le résultat spectral est central pour l'action BuP :
\[
S_{\rm BuP}[W]
=
\mathrm{Tr}\,L(W)^{-\beta[W]}
+
\sum_{ij}W_{ij}\kappa_{ij}[W]
+
\lambda\sum_{ij}W_{ij}d_{ij}^{2}
+
S_{\rm topo}[W].
\]

Si
\[
c_NL_N\to-\Delta_g,
\]
alors le terme spectral discret :
\[
\mathrm{Tr}\,L_N^{-\beta}
\]
peut être relié à une action spectrale continue :
\[
\mathrm{Tr}\,(-\Delta_g)^{-\beta}.
\]

Via la représentation de Mellin :
\[
\mathrm{Tr}\,L^{-\beta}
=
\frac{1}{\Gamma(\beta)}
\int_0^\infty
t^{\beta-1}
\mathrm{Tr}(e^{-tL})\,dt,
\]
et l'expansion de la trace de chaleur :
\[
\mathrm{Tr}(e^{-t\Delta_g})
\sim
(4\pi t)^{-D/2}
\sum_{n=0}^{\infty}
a_n(\Delta_g)t^n,
\]
le terme spectral contient des contributions géométriques continues :
\[
a_0\sim\int d^Dx\sqrt{|g|},
\]
et, selon la convention d'opérateur,
\[
a_1\ \text{ou}\ a_2
\sim
\int d^Dx\sqrt{|g|}\,R.
\]

Ainsi, la convergence spectrale étudiée dans Paper~16 est une condition
nécessaire pour que l'action BuP discrète possède une limite géométrique
continue.

\section{Limites}

Les résultats actuels sont positifs mais plusieurs limites doivent être
clairement énoncées.

Premièrement, \(c_N\) est encore ajusté numériquement. Une preuve complète
nécessite une dérivation analytique de ce facteur.

Deuxièmement, les tests sont effectués sur des géométries synthétiques
contrôlées. Il reste à étendre la méthode à des graphes d'intrication
quantiques où
\[
W_{ij}=I(i:j)
\]
n'est pas imposé comme un noyau gaussien.

Troisièmement, la convergence est testée uniquement sur les bas modes. Les
hauts modes sont plus sensibles à la discrétisation, au choix du graphe
\(k\)-NN et au bruit numérique.

Quatrièmement, les variétés testées sont compactes et sans bord. Les
conditions de bord devront être traitées séparément.

Cinquièmement, la convergence de valeurs propres n'implique pas à elle seule
la convergence complète des opérateurs. Elle doit être complétée par une
analyse des fonctions propres, des formes quadratiques ou des noyaux de
chaleur.

\section{Conclusion}

Paper~16 transforme la validation de dimension spectrale de Paper~15 en une
validation spectrale plus fine. Le Laplacien BuP rescalé
\[
L_N=\frac{D_N-W_N}{\epsilon_N}
\]
reconstruit le bas spectre du Laplacien de Laplace--Beltrami sur deux
géométries contrôlées.

Sur le cercle \(S^1\), les 12 premiers modes non nuls sont reconstruits avec
une erreur relative moyenne de :
\[
0.0051.
\]

Sur le tore plat \(T^2\), les 20 premiers modes non nuls sont reconstruits
avec une erreur relative moyenne de :
\[
0.0266.
\]

Ces résultats soutiennent la limite :
\[
c_NL_N
\longrightarrow
-\Delta_g.
\]

Le résultat n'est pas encore une preuve analytique. Il indique cependant que
la direction est correcte : le Laplacien d'intrication BuP, correctement
rescalé et normalisé, possède le bas spectre attendu d'un opérateur de
Laplace--Beltrami.

La prochaine étape consiste à dériver la normalisation \(c_N\), à préciser le
mode de convergence, puis à étendre le résultat des géométries contrôlées aux
graphes d'information mutuelle issus d'états quantiques.

\section*{Disponibilité du code et des données}

Les scripts et résultats de Paper~16 sont organisés dans le dossier :
\[
\texttt{papers/paper16\_spectral\_continuum\_limit/}.
\]

Les scripts principaux sont :
\[
\texttt{scripts/paper16\_circle\_spectrum\_convergence\_v1.py},
\]
\[
\texttt{scripts/paper16\_flat\_torus\_spectrum\_convergence\_v1.py},
\]
\[
\texttt{scripts/paper16\_build\_spectral\_summary\_v1.py}.
\]

Les résultats principaux sont stockés dans :
\[
\texttt{results/circle\_spectrum\_convergence\_v1/},
\]
\[
\texttt{results/flat\_torus\_spectrum\_convergence\_v1/},
\]
\[
\texttt{results/paper16\_spectral\_summary\_v1/}.
\]

\appendix

\section{Détails sur les erreurs relatives}

Pour chaque mode non nul \(k\), l'erreur relative est définie par :
\[
e_k
=
\frac{
\left|
c_N\lambda_k^{\rm graph}
-
\lambda_k^{\rm target}
\right|
}{
\lambda_k^{\rm target}
}.
\]

Les statistiques rapportées sont :
\[
\bar e
=
\frac{1}{K}
\sum_{k=1}^K e_k,
\]
\[
e_{\rm med}
=
\operatorname{median}(e_k),
\]
et
\[
e_{\max}
=
\max_k e_k.
\]

\section{Prochaines étapes}

Les prochaines étapes de Paper~16 sont :

\begin{enumerate}
    \item dériver analytiquement la normalisation \(c_N\) ;
    \item tester \(N=2048\) et \(N=4096\) avec méthodes sparse ;
    \item tester la convergence des fonctions propres ;
    \item comparer les formes quadratiques discrètes et continues ;
    \item tester une sphère avec spectre analytique
    \[
    \lambda_\ell=\ell(\ell+1)
    \]
    et dégénérescence \(2\ell+1\) ;
    \item remplacer les noyaux synthétiques par des graphes d'information
    mutuelle issus de véritables états quantiques ;
    \item relier le résultat à l'action spectrale de Paper~14 et à la limite
    Einstein de Paper~15.
\end{enumerate}

\end{document}
